% Figure 07 — Diagramme de séquence : Création / lecture de note
\documentclass[border=10pt]{standalone}
% Préambule commun pour tous les diagrammes TikZ standalone
% À inclure via % Préambule commun pour tous les diagrammes TikZ standalone
% À inclure via % Préambule commun pour tous les diagrammes TikZ standalone
% À inclure via \input{preambule.tex} après \documentclass{standalone}
\usepackage[utf8]{inputenc}
\usepackage[T1]{fontenc}
\usepackage[french]{babel}
\usepackage{lmodern}
\usepackage{tikz}
\usepackage{xcolor}

\usetikzlibrary{
  arrows.meta,
  positioning,
  shapes.geometric,
  shapes.symbols,
  shapes.misc,
  fit,
  calc,
  backgrounds,
  decorations.pathreplacing,
  decorations.markings,
  shadows.blur,
  patterns
}

% Palette de couleurs (cohérente sur tout le rapport)
\definecolor{awsorange}{HTML}{FF9900}
\definecolor{awsdark}{HTML}{232F3E}
\definecolor{awsblue}{HTML}{1A476F}
\definecolor{dockerblue}{HTML}{2496ED}
\definecolor{dockerdark}{HTML}{0DB7ED}
\definecolor{nginxgreen}{HTML}{009639}
\definecolor{flaskgray}{HTML}{3776AB}
\definecolor{pgblue}{HTML}{336791}
\definecolor{tlsgreen}{HTML}{2E8B57}
\definecolor{netgray}{HTML}{6C757D}
\definecolor{bgsoft}{HTML}{F5F7FA}
\definecolor{bordersoft}{HTML}{C8D1DC}
\definecolor{textdark}{HTML}{1F2937}
\definecolor{accentred}{HTML}{C0392B}
\definecolor{accentyellow}{HTML}{F1C40F}

% Styles génériques réutilisés dans plusieurs diagrammes
\tikzset{
  base/.style={
    rounded corners=4pt,
    draw=bordersoft,
    fill=white,
    text=textdark,
    align=center,
    inner sep=6pt,
    font=\small
  },
  zone/.style={
    base,
    fill=bgsoft,
    draw=netgray,
    dashed,
    inner sep=12pt
  },
  cloud/.style={
    base,
    fill=awsdark!5,
    draw=awsorange,
    line width=0.8pt,
    inner sep=14pt
  },
  service/.style={
    base,
    line width=0.6pt,
    minimum width=2.6cm,
    minimum height=1.1cm,
    drop shadow={shadow xshift=1pt, shadow yshift=-1pt, opacity=0.2}
  },
  legendebox/.style={
    base,
    fill=white,
    font=\footnotesize,
    inner sep=4pt
  },
  fleche/.style={
    -{Latex[length=2.5mm,width=2mm]},
    line width=0.6pt,
    color=netgray
  },
  fleche-bidir/.style={
    {Latex[length=2.5mm,width=2mm]}-{Latex[length=2.5mm,width=2mm]},
    line width=0.6pt,
    color=netgray
  },
  fleche-bold/.style={
    -{Latex[length=3mm,width=2.4mm]},
    line width=1pt
  },
  etiquette/.style={
    font=\scriptsize\sffamily,
    fill=white,
    inner sep=2pt,
    text=textdark
  }
}
 après \documentclass{standalone}
\usepackage[utf8]{inputenc}
\usepackage[T1]{fontenc}
\usepackage[french]{babel}
\usepackage{lmodern}
\usepackage{tikz}
\usepackage{xcolor}

\usetikzlibrary{
  arrows.meta,
  positioning,
  shapes.geometric,
  shapes.symbols,
  shapes.misc,
  fit,
  calc,
  backgrounds,
  decorations.pathreplacing,
  decorations.markings,
  shadows.blur,
  patterns
}

% Palette de couleurs (cohérente sur tout le rapport)
\definecolor{awsorange}{HTML}{FF9900}
\definecolor{awsdark}{HTML}{232F3E}
\definecolor{awsblue}{HTML}{1A476F}
\definecolor{dockerblue}{HTML}{2496ED}
\definecolor{dockerdark}{HTML}{0DB7ED}
\definecolor{nginxgreen}{HTML}{009639}
\definecolor{flaskgray}{HTML}{3776AB}
\definecolor{pgblue}{HTML}{336791}
\definecolor{tlsgreen}{HTML}{2E8B57}
\definecolor{netgray}{HTML}{6C757D}
\definecolor{bgsoft}{HTML}{F5F7FA}
\definecolor{bordersoft}{HTML}{C8D1DC}
\definecolor{textdark}{HTML}{1F2937}
\definecolor{accentred}{HTML}{C0392B}
\definecolor{accentyellow}{HTML}{F1C40F}

% Styles génériques réutilisés dans plusieurs diagrammes
\tikzset{
  base/.style={
    rounded corners=4pt,
    draw=bordersoft,
    fill=white,
    text=textdark,
    align=center,
    inner sep=6pt,
    font=\small
  },
  zone/.style={
    base,
    fill=bgsoft,
    draw=netgray,
    dashed,
    inner sep=12pt
  },
  cloud/.style={
    base,
    fill=awsdark!5,
    draw=awsorange,
    line width=0.8pt,
    inner sep=14pt
  },
  service/.style={
    base,
    line width=0.6pt,
    minimum width=2.6cm,
    minimum height=1.1cm,
    drop shadow={shadow xshift=1pt, shadow yshift=-1pt, opacity=0.2}
  },
  legendebox/.style={
    base,
    fill=white,
    font=\footnotesize,
    inner sep=4pt
  },
  fleche/.style={
    -{Latex[length=2.5mm,width=2mm]},
    line width=0.6pt,
    color=netgray
  },
  fleche-bidir/.style={
    {Latex[length=2.5mm,width=2mm]}-{Latex[length=2.5mm,width=2mm]},
    line width=0.6pt,
    color=netgray
  },
  fleche-bold/.style={
    -{Latex[length=3mm,width=2.4mm]},
    line width=1pt
  },
  etiquette/.style={
    font=\scriptsize\sffamily,
    fill=white,
    inner sep=2pt,
    text=textdark
  }
}
 après \documentclass{standalone}
\usepackage[utf8]{inputenc}
\usepackage[T1]{fontenc}
\usepackage[french]{babel}
\usepackage{lmodern}
\usepackage{tikz}
\usepackage{xcolor}

\usetikzlibrary{
  arrows.meta,
  positioning,
  shapes.geometric,
  shapes.symbols,
  shapes.misc,
  fit,
  calc,
  backgrounds,
  decorations.pathreplacing,
  decorations.markings,
  shadows.blur,
  patterns
}

% Palette de couleurs (cohérente sur tout le rapport)
\definecolor{awsorange}{HTML}{FF9900}
\definecolor{awsdark}{HTML}{232F3E}
\definecolor{awsblue}{HTML}{1A476F}
\definecolor{dockerblue}{HTML}{2496ED}
\definecolor{dockerdark}{HTML}{0DB7ED}
\definecolor{nginxgreen}{HTML}{009639}
\definecolor{flaskgray}{HTML}{3776AB}
\definecolor{pgblue}{HTML}{336791}
\definecolor{tlsgreen}{HTML}{2E8B57}
\definecolor{netgray}{HTML}{6C757D}
\definecolor{bgsoft}{HTML}{F5F7FA}
\definecolor{bordersoft}{HTML}{C8D1DC}
\definecolor{textdark}{HTML}{1F2937}
\definecolor{accentred}{HTML}{C0392B}
\definecolor{accentyellow}{HTML}{F1C40F}

% Styles génériques réutilisés dans plusieurs diagrammes
\tikzset{
  base/.style={
    rounded corners=4pt,
    draw=bordersoft,
    fill=white,
    text=textdark,
    align=center,
    inner sep=6pt,
    font=\small
  },
  zone/.style={
    base,
    fill=bgsoft,
    draw=netgray,
    dashed,
    inner sep=12pt
  },
  cloud/.style={
    base,
    fill=awsdark!5,
    draw=awsorange,
    line width=0.8pt,
    inner sep=14pt
  },
  service/.style={
    base,
    line width=0.6pt,
    minimum width=2.6cm,
    minimum height=1.1cm,
    drop shadow={shadow xshift=1pt, shadow yshift=-1pt, opacity=0.2}
  },
  legendebox/.style={
    base,
    fill=white,
    font=\footnotesize,
    inner sep=4pt
  },
  fleche/.style={
    -{Latex[length=2.5mm,width=2mm]},
    line width=0.6pt,
    color=netgray
  },
  fleche-bidir/.style={
    {Latex[length=2.5mm,width=2mm]}-{Latex[length=2.5mm,width=2mm]},
    line width=0.6pt,
    color=netgray
  },
  fleche-bold/.style={
    -{Latex[length=3mm,width=2.4mm]},
    line width=1pt
  },
  etiquette/.style={
    font=\scriptsize\sffamily,
    fill=white,
    inner sep=2pt,
    text=textdark
  }
}


\begin{document}
\begin{tikzpicture}[
  acteur/.style={
    base, fill=awsblue!10, draw=awsblue, line width=0.6pt,
    minimum width=2.4cm, minimum height=0.8cm,
    font=\small\bfseries
  },
  msg/.style={-{Latex[length=2mm]}, line width=0.6pt, color=textdark},
  msg-ret/.style={-{Latex[length=2mm]}, dashed, line width=0.6pt, color=netgray},
  activation/.style={fill=awsorange!40, draw=awsorange, line width=0.4pt},
  note/.style={
    base, fill=accentyellow!15, draw=accentyellow!80!black,
    font=\scriptsize, align=left, inner sep=3pt
  }
]

  \def\xUser{0}
  \def\xBrowser{3.2}
  \def\xNginx{6.4}
  \def\xFlask{9.6}
  \def\xDB{12.8}

  \node[acteur] (user)    at (\xUser,    0) {Utilisateur};
  \node[acteur] (browser) at (\xBrowser, 0) {Navigateur};
  \node[acteur] (nginx)   at (\xNginx,   0) {Nginx};
  \node[acteur] (flask)   at (\xFlask,   0) {Flask};
  \node[acteur] (db)      at (\xDB,      0) {PostgreSQL};

  \def\bottom{-13.0}
  \foreach \x in {\xUser, \xBrowser, \xNginx, \xFlask, \xDB} {
    \draw[dashed, color=netgray] (\x, -0.4) -- (\x, \bottom);
  }

  % --- Phase 1 : affichage formulaire ---
  \draw[msg] (\xUser, -1.0) -- node[etiquette, above]{1. Clic « Nouvelle note »}
    (\xBrowser, -1.0);
  \draw[msg] (\xBrowser, -1.7) -- node[etiquette, above]{2. GET \texttt{/notes/create}}
    (\xNginx, -1.7);
  \fill[activation] (\xNginx-0.1, -1.7) rectangle (\xNginx+0.1, -3.5);

  \draw[msg] (\xNginx, -2.2) -- node[etiquette, above]{3. proxy\_pass}
    (\xFlask, -2.2);
  \fill[activation] (\xFlask-0.1, -2.2) rectangle (\xFlask+0.1, -3.0);

  \draw[msg-ret] (\xFlask, -2.8) -- node[etiquette, above, text=tlsgreen]{4. HTML formulaire (CSRF token inclus)}
    (\xNginx, -2.8);
  \draw[msg-ret] (\xNginx, -3.3) -- node[etiquette, above]{5. réponse HTTPS}
    (\xBrowser, -3.3);

  % --- Phase 2 : soumission ---
  \draw[msg] (\xUser, -4.5) -- node[etiquette, above]{6. Saisit titre + contenu, valide}
    (\xBrowser, -4.5);
  \draw[msg] (\xBrowser, -5.2) -- node[etiquette, above]{7. POST \texttt{/notes/create} (titre, contenu, csrf\_token)}
    (\xNginx, -5.2);

  \fill[activation] (\xNginx-0.1, -5.2) rectangle (\xNginx+0.1, -11.0);
  \draw[msg] (\xNginx, -5.8) -- node[etiquette, above]{8. proxy\_pass}
    (\xFlask, -5.8);
  \fill[activation] (\xFlask-0.1, -5.8) rectangle (\xFlask+0.1, -10.4);

  % --- Étape 9 : @login_required + CSRF ---
  \draw[msg] (\xFlask, -6.4) to[bend left=90, looseness=4]
    node[etiquette, right, xshift=2mm]{9. \texttt{@login\_required} \(+\) CSRF}
    (\xFlask, -7.0);

  % --- Étape 10 : INSERT ---
  \draw[msg, color=accentred, line width=0.8pt]
    (\xFlask, -7.8) -- node[etiquette, above, text=accentred]
      {10. \texttt{INSERT INTO notes (title, content, user\_id) VALUES (\$1, \$2, current\_user.id)}}
    (\xDB, -7.8);
  \fill[activation] (\xDB-0.1, -7.8) rectangle (\xDB+0.1, -8.6);

  % --- Étape 11 : retour id ---
  \draw[msg-ret] (\xDB, -8.4) -- node[etiquette, above, text=netgray]{11. \texttt{RETURNING id}}
    (\xFlask, -8.4);

  % --- Étape 12 : redirection ---
  \draw[msg-ret] (\xFlask, -9.6) -- node[etiquette, above, text=tlsgreen]{12. 302 \(\to\) \texttt{/notes}}
    (\xNginx, -9.6);
  \draw[msg-ret] (\xNginx, -10.6) -- node[etiquette, above, text=tlsgreen]{13. réponse HTTPS}
    (\xBrowser, -10.6);
  \draw[msg-ret] (\xBrowser, -11.3) -- node[etiquette, above]{14. Nouvelle note affichée}
    (\xUser, -11.3);

  % --- Note d'isolation ---
  \node[note, anchor=north west, text width=12cm]
    at (0.5, -12.0)
    {\textbf{Isolation par utilisateur :} toutes les requêtes de lecture filtrent par
     \texttt{WHERE user\_id = current\_user.id}. Aucun utilisateur ne peut accéder
     aux notes d'un autre, même en forçant un identifiant dans l'URL (IDOR bloqué).};

\end{tikzpicture}
\end{document}
